\documentclass[11pt]{article}
\usepackage[margin=1in]{geometry}
\usepackage[hidelinks]{hyperref}
\usepackage{parskip}

\title{Living-Documentation Cleanup}
\author{Haoze Wu (with Codex)}
\date{2026-08-20}

\begin{document}
\maketitle

\section{Objective and outcome}

The objective was to remove obsolete and redundant material from the mutable
project documentation without weakening its contracts, provenance, or resume
value.  The cleanup is complete.  \texttt{PROGRESS.md} fell from 411 to 78
lines, and \texttt{docs/scheduling-parity-ledger.md} fell from 597 to 457
lines.  Across the two files, the change removes 562 lines and adds 89.

No solver source, benchmark artifact, frozen design snapshot, historical
report, or documentation checker changed.  No whole documentation file was
deleted.

\section{Decisions and rationale}

\begin{itemize}
\item \texttt{PROGRESS.md} is now only the current resume surface.  Old
  rank-by-rank investigation diaries, the duplicated pre-Phase-4.5 debug
  contract, and the duplicated later-phase roadmap were removed because their
  durable owners are the development plan, frozen benchmark reports, design
  snapshots, and the scheduling ledger.
\item The progress file retains the current Phase-4 exit state, the six closed
  Tier-D ranks, the fourteen open ranks, the latest soundness evidence, open
  gates, immediate work, and the machine-parsed scheduling markers.
\item Only the mutable N4 verification tail of the scheduling ledger was
  compacted.  Superseded ``validation pending'' chronology and rejected
  candidate diaries were replaced by the accepted final disposition, checked
  job provenance, later proof-soundness evidence, and explicit reopening
  conditions.
\item Frozen design documents and reports were retained because they preserve
  distinct evidence and historical provenance.  All parser anchors, ledger
  row identifiers, option-accounting sections, and policy boundaries were
  retained.
\end{itemize}

\section{Validation and provenance}

This is a documentation-only authority cleanup; the solver binary and its
behavior are unchanged, so solver integration gates were not rerun.  Measured
results from the full documentation checker suite are: zero citation errors,
zero section-reference errors (80/80 references resolved), zero progress
errors or warnings, 28/28 inventory entries accounted for, 418/418 reference
options accounted for, and no failures in proof-rule synchronization,
configuration wiring, default-strategy configuration, effort specifications,
statistic identifiers, comment hygiene, or rule/policy boundaries.

The size checker also passes.  It reports two non-failing reserve warnings:
\texttt{docs/calculus.md} is 635/680 lines and
\texttt{docs/implementation-notes.md} is 559/560 lines.  Those files were not
edited because this audit found no safely removable current content in them.

\section{Limitations and risks}

The cleanup deliberately does not rewrite immutable evidence or attempt a
second semantic consolidation of the calculus and implementation notes.
Historical debugging detail removed from the live progress surface remains in
the cited frozen campaign reports and repository history.  The main ongoing
risk is future duplication between the progress file and its authoritative
owners; the new ownership preamble and existing size checker make that drift
visible.

An independent Claude review was requested under the normal documentation
workflow, but the external-service handoff was unavailable under the active
data-export policy.  The user explicitly approved proceeding with the local
audit and diagnosis.

\section{Next actions}

Resume the development plan from the concise \texttt{PROGRESS.md} surface:
keep rank 22 parked for now, investigate another open Tier-D rank under the
pre-Phase-4.5 contract, and leave N4 closed unless a documented reopening
condition occurs.

\end{document}
